\documentclass[10pt,a4paper]{article}
\usepackage[utf8]{inputenc}
\usepackage[T1]{fontenc}
\usepackage{lmodern}
\usepackage[margin=2.2cm]{geometry}
\usepackage{amsmath,amssymb}
\usepackage{graphicx}
\usepackage{booktabs}
\usepackage{caption}
\usepackage{float}
\usepackage{hyperref}
\captionsetup{font=small,labelfont=bf}
\input{values}

\title{From a binary line map to charge states $(n_1,n_2)$\\[2pt]
\large method, step by step, and results}
\author{\small Generated by \texttt{docs/charge\_states/make\_note.py}\\
\small run: \texttt{4-5-6-7-8\_rays\_40-50-60\_points\_500\_samples}}
\date{}

\begin{document}
\maketitle

\noindent
This note starts where the U-Net stops: a binary map of transition lines.
It describes how the charge state $(n_1,n_2)$ of every point of the gate
plane is read off that map, and how well this works.  The same function,
\texttt{label\_states}, is applied in exactly the same way to two inputs:

\begin{itemize}
\item the full CSD line map $Y$, which gives the ceiling of the method;
\item the U-Net map $B = P > t$, with $P$ the U-Net probability and $t$
  the threshold stored in each \texttt{unet.pt} (chosen on the validation
  devices during training, not re-tuned here).
\end{itemize}

\noindent
Every number below is written by the program from the files it read.
Figures~\ref{fig:regions}--\ref{fig:scoring} follow one test device,
\vDevice{} (\vNStatesDev{} charge states in the window), through every step,
on $Y$ (top rows) and on the $8\times60$ U-Net with \vCoverage\,\% of the
grid measured (bottom rows).  The code is
\texttt{src/csdrecon/study/charge\_states.py}; the study is
\texttt{scripts/run\_16\_charge\_states.py}.

\section{Conventions}

The map is $H\times W = 100\times100$ pixels; row $=V_\mathrm{P2}$ with row~0
lowest, column $=V_\mathrm{P1}$.  The pixel size $(\Delta x,\Delta y)$ in mV
comes from \texttt{voltage\_window} in \texttt{device.json}.  The carriers are
holes, so the occupation grows towards \emph{negative} gate voltage: $n_1$
grows to lower $V_\mathrm{P1}$, $n_2$ to lower $V_\mathrm{P2}$, and $(0,0)$ lies
towards the top right.  This was checked on \texttt{charge\_states.npy}, not
assumed: on all \vNTest{} test devices $n_1$ never grows towards
$+V_\mathrm{P1}$, $n_2$ never towards $+V_\mathrm{P2}$, and the top-right corner
is $(0,0)$ on \vCornerOO/\vNTest.

\section{The method}

\subsection*{Step 1 --- close the frame}
The simulator's line maps stop one pixel short of the image edge; two
regions then join around the end of a line.  The line pixels of the second
row/column are copied onto the outermost one (\texttt{close\_frame}).  Without
this step, labelling the full CSD reached only 39\,\% accuracy on the first
training devices checked; with it, 100\,\%.  It is applied identically to
both inputs (red pixels in Fig.~\ref{fig:regions}b,f).

\subsection*{Step 2 --- regions}
Connected components of the non-line pixels $\lnot B$, with 4-connectivity
(two pixels touch only through an edge, so a one-pixel diagonal line still
separates them).  Components smaller than \vMinRegion{} pixels are crumbs
between close lines and are counted as line (Fig.~\ref{fig:regions}c,g).
Each remaining component is one candidate charge state.

\subsection*{Step 3 --- fill the lines}
Every line pixel takes the label of the nearest region (Euclidean distance
transform with indices).  Every pixel now belongs to a region
(Fig.~\ref{fig:regions}d,h).

\begin{figure}[H]
\centering
\includegraphics[width=\textwidth]{figures/step1_regions.pdf}
\caption{Steps 1--3 on \vDevice.  Top: full CSD, \vNRegFull{} regions.
Bottom: U-Net $8\times60$, \vNRegUnet{} regions.  The probability map (e)
does show the lines, but the binary map (f) has gaps in them, and a gap lets
two states fall into one region.}
\label{fig:regions}
\end{figure}

\subsection*{Step 4 --- borders}
Two regions are neighbours if a pixel of one is a 4-neighbour of a pixel of
the other.  Each such pair of pixels is a \emph{contact}; the border length
is the number of contacts.  The border's direction is the principal axis of
its contact midpoints $(x_k,y_k)$ in mV, centred on their mean:
\begin{equation}
\theta = \tfrac12\,\operatorname{atan2}\!\bigl(2S_{xy},\;S_{xx}-S_{yy}\bigr),
\qquad S_{xy}=\textstyle\sum_k x_k y_k,\ \dots
\end{equation}
taken in $[0^\circ,180^\circ)$.  The border normal is the mean, over the
contacts, of the step (in mV) from the pixel in region $b$ to the pixel in
region $a$; it points into $a$.  Borders with fewer than \vMinBorder{}
contacts get no direction and are not used.

\subsection*{Step 5 --- family of each border}
\begin{itemize}
\item positive slope ($0^\circ<\theta<90^\circ$) $\to$ \textbf{interdot}:
  the charge moves between the dots;
\item negative slope $\to$ \textbf{dot 1} or \textbf{dot 2}.  The negative
  angles of the whole device are split in two by Otsu's threshold, weighted
  by border length.  The steeper group (closer to $90^\circ$) is dot~1, the
  shallower is dot~2.  An exactly horizontal border counts as the shallowest
  angle, an exactly vertical one as the steepest.  A device with a single
  negative-slope border is split at $135^\circ$.
\end{itemize}
Figure~\ref{fig:borders}c,f shows the length-weighted angle histogram, the
Otsu cut and the true angles from the capacitances.

\subsection*{Step 6 --- which side gets the extra charge}
From the border normal $(n_x,n_y)$, pointing from $b$ into $a$:

\begin{center}
\begin{tabular}{lll}
\toprule
family & $a$ gains when & $(n_1,n_2)_a-(n_1,n_2)_b$ \\
\midrule
dot 1    & $n_x<0$ ($a$ at lower $V_\mathrm{P1}$) & $\pm(1,0)$ \\
dot 2    & $n_y<0$ ($a$ at lower $V_\mathrm{P2}$) & $\pm(0,1)$ \\
interdot & $n_x<0$ ($a$ on the lower-$V_\mathrm{P1}$ side) & $\pm(1,-1)$ \\
\bottomrule
\end{tabular}
\end{center}

\noindent
with $+$ when the condition holds and $-$ otherwise
(Fig.~\ref{fig:borders}b,e: each arrow points into the side that gains).

\begin{figure}[H]
\centering
\includegraphics[width=\textwidth]{figures/step2_borders.pdf}
\caption{Steps 4--6.  (a,d) Border contacts coloured by family.
(b,e) The step across each border.  (c,f) Border angles, length-weighted, with
the Otsu cut (dashed) and the true angles from the capacitance matrices
(solid).  On the full CSD the two dot families fall on their true angles; on
the U-Net map few borders are left to classify.}
\label{fig:borders}
\end{figure}

\subsection*{Step 7 --- propagate the labels}
The regions are the nodes of a graph and the borders with a direction are
its edges.  Its maximum spanning tree by border length is kept, because long
borders give the most reliable angle and normal.  The region holding the
top-right pixel (highest $V_\mathrm{P1}$, $V_\mathrm{P2}$) is the root and gets
$(0,0)$.  The tree is walked breadth-first, adding the step of each border.
A region the tree does not reach stays unlabelled and is scored as wrong.

\subsection*{Step 8 --- consistency}
Every border \emph{not} in the tree gives a second route between two regions
that already have labels.  The \emph{consistency} of a device is the fraction
of these borders whose step agrees with the labels.  It needs no ground
truth: 1.0 means the line map is self-consistent, and a lower value means
that lines are missing or invented (Fig.~\ref{fig:tree}a,d: green and red
dashed edges).

\begin{figure}[H]
\centering
\includegraphics[width=\textwidth]{figures/step3_tree_labels.pdf}
\caption{Steps 7--8.  (a,d) Region graph: maximum spanning tree (solid) and
the borders left out of it (dashed), coloured by whether they agree with the
labels.  (b,e) The labels.  (c,f) The truth.  Full CSD: \vAccDevFull\,\%
correct.  U-Net: \vAccDevUnet\,\%, \vRecDevUnet\,\% of the states recovered,
consistency \vConDevUnet.}
\label{fig:tree}
\end{figure}

\section{Scoring}

All metrics are computed per device and then averaged over devices.

\begin{itemize}
\item \textbf{Scored pixels.}  Pixels within 1\,px of a true line are left
  out: their state is ambiguous at this resolution, and this matches the
  1-pixel tolerance of F1@1.  On \vDevice, \vScoredDev\,\% of the pixels are
  scored.
\item \textbf{Relative pixel accuracy.}  The best constant shift
  $(\Delta n_1,\Delta n_2)$ per device is applied, then the fraction of
  scored pixels with the right $(n_1,n_2)$ is taken.  This tests the
  \emph{pattern}.
\item \textbf{Absolute pixel accuracy.}  The same with no shift.  It is
  only taken on devices whose true corner is $(0,0)$, which here is all of
  them.
\item \textbf{Total-charge accuracy.}  The same for $n_1+n_2$, with its own
  best shift.
\item \textbf{State recall.}  The fraction of the true states whose scored
  pixels are mostly ($>50\,\%$) labelled right.  \textbf{State precision}:
  the fraction of predicted regions whose scored pixels are mostly the state
  they were labelled.
\item \textbf{(1,1) IoU and centre error.}  For the state used for a qubit,
  computed with the shifted labels.  The centre error is the distance in mV
  between the true and the predicted (1,1) centroids.
\end{itemize}

\paragraph{A baseline that must be beaten.}  The $(0,0)$ region is large:
labelling \emph{every} pixel $(0,0)$ already gives
\vBaseMean\,$\pm$\,\vBaseStd\,\% pixel accuracy on the test devices
(\vBaseDev\,\% on \vDevice, Fig.~\ref{fig:scoring}c).  Pixel accuracy alone
therefore overstates the result; state recall and the (1,1) IoU do not have
this problem.

\begin{figure}[H]
\centering
\includegraphics[width=0.78\textwidth]{figures/step4_scoring.pdf}
\caption{Scoring \vDevice.  (a) Scored pixels (white: within 1\,px of a true
line).  (b) U-Net errors; orange: $n_1$ wrong, blue: $n_2$ wrong, black:
both.  (c) The trivial ``all (0,0)'' labelling scores almost the same as the
U-Net on this device.}
\label{fig:scoring}
\end{figure}

\section{Results}

\subsection*{How hard the task is}
Capacitances are drawn at random and the window is randomly offset, so the
devices differ a lot (Fig.~\ref{fig:variation}).  The test devices have
\vNStMin--\vNStMax{} charge states in the window (median \vNStMed), and a
median region size of \vRegMin--\vRegMax{} pixels.

\begin{figure}[H]
\centering
\includegraphics[width=\textwidth]{figures/fig1_device_variation.pdf}
\caption{Device variation over the \vNTest{} test devices, and six examples
spread over the number of states.}
\label{fig:variation}
\end{figure}

\subsection*{Tables}

\begin{table}[H]
\centering
\caption{Test split (\vNTest{} devices, held out).  Each cell: mean $\pm$ std
over devices, median in brackets.}
\resizebox{\textwidth}{!}{\input{summary.tex}}
\end{table}

\begin{table}[H]
\centering
\caption{Validation split (75 devices).  These devices chose each U-Net's
threshold, so they are \emph{not} held out.}
\resizebox{\textwidth}{!}{\input{summary_validation.tex}}
\end{table}

\subsection*{Summary of the numbers}
On the full CSD the method labels \vFullAcc\,\% of the scored plane
correctly and recovers \vFullRec\,\% of the charge states, so the labelling
step itself is not the limit.  From the U-Net maps the relative pixel
accuracy is \vLoAcc\,\% at \vLoCov\,\% coverage ($4\times40$) and \vMidAcc\,\%
(median \vMidAccMed\,\%) at \vMidCov\,\% ($8\times60$, the best U-Net of this
sweep).  This is to be read against the \vBaseMean\,\% of the ``all (0,0)''
baseline.  The U-Nets recover \vLoRec--\vMidRec\,\% of the charge states,
against \vFullRec\,\% from the full CSD.  The (1,1) region has a mean IoU of
\vMidIoU{} (median \vMidIoUMed) for $8\times60$, and the line maps'
consistency is \vMidCon, against 1.00 for the full CSD.

\begin{figure}[H]
\centering
\includegraphics[width=\textwidth]{figures/fig2_labelling_examples.pdf}
\caption{Good, typical and bad test device for the $8\times60$ U-Net, by
relative pixel accuracy.}
\label{fig:examples}
\end{figure}

\begin{figure}[H]
\centering
\includegraphics[width=0.9\textwidth]{figures/fig3_accuracy_vs_coverage.pdf}
\caption{Relative pixel accuracy and state recall against the fraction of
the grid measured, one point per U-Net; the full CSD at 100\,\%.}
\label{fig:coverage}
\end{figure}

\begin{figure}[H]
\centering
\includegraphics[width=\textwidth]{figures/fig4_failures.pdf}
\caption{What goes wrong for the $8\times60$ U-Net.  Left: the label error
$(\Delta n_1,\Delta n_2)$ of the wrong pixels, all test devices.  Middle,
right: per-device accuracy against the number of states in the window and
against the interdot capacitance $d_1d_2$.}
\label{fig:failures}
\end{figure}

\section{Limitations}
\begin{itemize}
\item \textbf{One gap merges two states.}  A break in a predicted line joins
  two regions.  The labels beyond that border are then all wrong by one
  step, and the error spreads through the tree to every region behind it.
  This is the dominant failure (Figs.~\ref{fig:regions}f--h and
  \ref{fig:examples}).
\item \textbf{The corner assumption.}  Absolute labels assume the top-right
  region is $(0,0)$.  This holds for every device in this pool, but it is a
  property of the simulated windows, not of the method.
\item \textbf{Degenerate borders.}  A border whose contacts point in
  opposite directions (for example one that wraps round a small region) has
  a zero mean normal, and its side is then decided by the sign convention
  alone.  This occurs on \vDevice{} for the U-Net map.
\end{itemize}

\end{document}
